\section{Déploiement Opérationnel et Validation des Services}

%------------------------------------------------------------
% Slide 11 : Services Réseau Fondamentaux (DHCP & DNS)
%------------------------------------------------------------
\begin{frame}{Services Réseau Fondamentaux : DHCP et DNS}
    Le serveur principal « Tout-en-un » (\texttt{10.0.20.10}) orchestre l'infrastructure logique.

    \vspace{0.2cm}
    \begin{columns}[T]
        \begin{column}{0.5\textwidth}
            \begin{block}{isc-dhcp-server}
                \begin{itemize}
                    \item Attribution automatique des baux IP aux stations de recherche.
                    \item Cycle \textbf{DORA} complet émulé et validé.
                    \item Injection automatique des passerelles et du DNS local.
                \end{itemize}
            \end{block}
        \end{column}
        
        \begin{column}{0.5\textwidth}
            \begin{block}{BIND9 (Serveur DNS)}
                \begin{itemize}
                    \item Résolution de noms locale pour les serveurs de l'institut.
                    \item Configuration de redirecteurs (\textit{forwarders}) vers le WAN.
                    \item Sécurisation des requêtes récursives.
                \end{itemize}
            \end{block}
        \end{column}
    \end{columns}
\end{frame}


\begin{frame}{Validation du Statut des Services Applicatifs (VLAN 20)}
    Pour répondre aux besoins des 80 chercheurs, l'infrastructure s'appuie sur trois services clés fonctionnels et sécurisés : le transfert de données (\texttt{vsftpd}), l'attribution dynamique d'IP (\texttt{dhcpd}) et la résolution de noms (\texttt{named}).

    \begin{columns}[T]
        \begin{column}{0.6\textwidth}
            \begin{figure}
                \centering
                \includegraphics[width=1\textwidth]{../rapport/figures/test2_services.png}
                \caption{Statut opérationnel des daemons systémd}
            \end{figure}
        \end{column}
        
        \begin{column}{0.4\textwidth}
            \textbf{Services Validés (Active/Running):}
            \begin{itemize}
                \item \textbf{\texttt{vsftpd.service}} : Serveur FTP actif pour le transfert sécurisé des chercheurs.
                \item \textbf{\texttt{isc-dhcp-server}} : Traitement des requêtes \texttt{DHCPREQUEST} et envois des \texttt{DHCPACK} pour les clients (ex: PC-Client-21).
                \item \textbf{\texttt{named.service}} : Serveur DNS BIND opérationnel avec isolation des processus.
            \end{itemize}
        \end{column}
    \end{columns}
\end{frame}


%------------------------------------------------------------
% Slide 12 : Preuve Technique - Attribution DHCP
%------------------------------------------------------------
\begin{frame}{Preuve Technique : Attribution Dynamique DHCP}
    \begin{columns}[T]
        \begin{column}{0.4\textwidth}
            \textbf{Validation du Cycle DORA :}
            \begin{itemize}
                \item Démarrage d'un nœud client sur la maquette.
                \item Capture des logs système mettant en évidence les requêtes et acquittements.
                \item Attribution réussie des paramètres réseau sans conflit.
            \end{itemize}
        \end{column}
        
        \begin{column}{0.6\textwidth}
            \begin{figure}
                \centering
                % Intégration de votre capture test_attribution_dhcp.png
                \includegraphics[width=0.98\textwidth]{figures/test_attribution_dhcp.png}
                \caption{Capture d'écran du processus d'attribution DHCP}
            \end{figure}
        \end{column}
    \end{columns}
\end{frame}

%------------------------------------------------------------
% Slide 13 : Preuve Technique - Routage & Connectivité
%------------------------------------------------------------
\begin{frame}{Preuve Technique : Connectivité et Routage Inter-VLAN}
    \begin{columns}[T]
        \begin{column}{0.6\textwidth}
            \begin{figure}
                \centering
                % Intégration de votre capture vision_globale.jpeg
                % \includegraphics[width=0.98\textwidth]{figures/test_.jpeg}
                \caption{Validation des pings croisés à travers le cœur de réseau}
            \end{figure}
        \end{column}
        
        \begin{column}{0.4\textwidth}
            \textbf{Validation du Cœur de Réseau :}
            \begin{itemize}
                \item Test de connectivité par commandes \texttt{ping} croisées.
                \item Transit réussi des paquets ICMP entre des sous-réseaux distincts.
                \item Accès validé vers le serveur de production (\texttt{VLAN 20}) et le serveur de sauvegarde (\texttt{VLAN 60}).
            \end{itemize}
        \end{column}
    \end{columns}
\end{frame}

%------------------------------------------------------------
% Slide 14 : Plateforme de Stockage et Partage (FTPS)
%------------------------------------------------------------



%------------------------------------------------------------
% Slide 15 : Preuve Technique - Transfert Clinique
%------------------------------------------------------------
\begin{frame}{Preuve Technique : Validation du Transfert FTP }

    \begin{columns}[T]
        \begin{column}{0.4\textwidth}
            \textbf{Simulation en conditions réelles :}
            \begin{itemize}
                \item Connexion établie depuis un client distant  \texttt{ftp}.
                \item Authentification et établissement du canal de données.
                \item Commande \texttt{put} exécutée avec succès sans aucune perte de paquets.
            \end{itemize}
        \end{column}
        
        \begin{column}{0.6\textwidth}
            \begin{figure}
                \centering
                % Intégration de votre capture test8trqnsfert8ftp.png
                \includegraphics[width=0.95\textwidth]{figures/test8trqnsfert8ftp.png}
                \caption{Flux de connexion et dépôt de fichier validés}
            \end{figure}
        \end{column}
    \end{columns}
\end{frame}